% =============================================================================
%  Notas del Profesor — Sesión 23: Argumentos de consola
%  Programación Avanzada — Universidad Panamericana
%  Compilar: pdflatex sesion23_argumentos_consola_profesor.tex
% =============================================================================
\documentclass[12pt, oneside]{article}
\usepackage[profesor]{progavanzada}
\usepackage{array}

\begin{document}

\portada{Sesión 23}{Argumentos de consola}{2026}

\tableofcontents
\newpage

% =============================================================================
\section{Objetivos de la sesión}
% =============================================================================

Al finalizar la sesión el alumno será capaz de:

\begin{enumerate}
  \item Explicar la firma completa de \cpp{main} con \cpp{argc} y
        \cpp{argv}, y describir el valor de cada parámetro para una
        ejecución concreta.
  \item Escribir programas que reciban argumentos de consola y los
        procesen con un ciclo sobre \cpp{argv}.
  \item Verificar correctamente el número de argumentos e imprimir un
        mensaje de uso estándar en \cpp{cerr}.
  \item Convertir argumentos de cadena a tipos numéricos con
        \cpp{atoi}, \cpp{atof} y \cpp{stoi}.
  \item Resolver la actividad del teclado numérico usando \cpp{argc},
        \cpp{argv} y una tabla de mapeo letra--tecla.
\end{enumerate}

\begin{notaprofesor}[Duración y distribución (100 min)]
  \begin{tabular}{rl}
    10 min & Gancho histórico + preguntas de arranque \\
    10 min & \cpp{main} con parámetros: \cpp{argc} y \cpp{argv} \\
    20 min & Demo en vivo: cuatro programas (ver sección~4) \\
     5 min & Tabla de conversión a números \\
    55 min & Actividad: teclado numérico \\
  \end{tabular}
  \par\smallskip
  La actividad es la parte más rica de la sesión.  Los primeros
  35~minutos se dedican a implementar; los 20~restantes a discutir
  soluciones y ver el diseño de la tabla de mapeo.
\end{notaprofesor}

% =============================================================================
\section{Prerrequisitos}
% =============================================================================

\begin{itemize}
  \item Ciclos \cpp{for} y arreglos (sesiones 7--8).
  \item Cadenas estilo C: \cpp{char*}, acceso por índice, terminador
        \texttt{'\textbackslash{}0'} (sesiones 3 y 9).
  \item Tabla ASCII y operaciones con \cpp{char} (sesión 3).
  \item Compilación y ejecución básica desde terminal con \texttt{g++}.
\end{itemize}

\begin{notaprofesor}
  Esta sesión no requiere archivos, structs ni memoria dinámica.  Es
  conceptualmente ligera: los alumnos ya conocen arreglos y ya usan la
  terminal.  El reto es que \cpp{argv} combina dos conceptos que dominan
  por separado (arreglos y cadenas) en un solo objeto que nunca habían
  visto así.
  \par\smallskip
  Señal de alerta: si algún alumno no recuerda que \cpp{char*} es un
  puntero a una cadena terminada en \texttt{'\textbackslash{}0'}, tomar
  3~minutos para mostrar que \cpp{argv[1][0]} es el primer carácter del
  primer argumento, antes de avanzar.  El ejercicio~4 (contador) lo
  necesita directamente.
\end{notaprofesor}

% =============================================================================
\section{Arranque y preguntas de sondeo}
% =============================================================================

\begin{notaprofesor}[Preguntas de arranque (10 min)]

  \textbf{1.\ ``¿Cómo le pasan instrucciones a \texttt{g++} cuando
  compilan?''}\\
  Respuesta esperada: escribiendo opciones en la línea de comandos
  (\texttt{-o}, \texttt{-Wall}, nombre del archivo).  Seguir con:
  ``eso significa que \texttt{g++} recibe esos textos de alguna manera
  dentro de su \cpp{main}.  Hoy van a aprender exactamente cómo.''

  \medskip
  \textbf{2.\ ``¿Cuántos argumentos recibe \texttt{g++} cuando
  ejecutan \texttt{g++ -O2 -o hola hola.cpp}?''}\\
  Respuesta esperada (varios intentos): la respuesta es~4.  El nombre
  del programa (\texttt{g++}) se cuenta aparte y es \cpp{argv[0]}.
  Este pequeño ``¿qué tal 4 o 5?'' crea la tensión cognitiva perfecta
  para explicar \cpp{argc}.

  \medskip
  \textbf{3.\ ``¿Qué diferencia hay entre \texttt{./programa} y
  \texttt{./programa hola mundo}?''}\\
  Respuesta esperada: en el segundo caso el programa recibe dos
  palabras extra.  Seguir con: ``¿cómo las recibe?  ¿Por dónde entran?
  ¿Desde \cpp{cin}?''  Respuesta esperada: no, \cpp{cin} lee del
  teclado \textit{durante} la ejecución; los argumentos de consola
  se entregan \textit{antes} de que el programa comience.

\end{notaprofesor}

% =============================================================================
\section{\texttt{main} con parámetros}
% =============================================================================

\begin{notaprofesor}[Cómo presentar \texttt{argc}/\texttt{argv} (10 min)]
  Escribir en la pizarra o proyectar los dos \cpp{main} uno junto al otro:
  \begin{verbatim}
int main()                           // el que siempre han usado
int main(int argc, char* argv[])     // la versión completa
  \end{verbatim}
  Preguntar: ``¿qué nuevo hay aquí?''  Identificar \cpp{int argc} y
  \cpp{char* argv[]}.  Explicar que \cpp{argv[]} es un arreglo de
  \cpp{char*}: cada elemento apunta a una cadena de texto terminada en
  \texttt{'\textbackslash{}0'}.

  \medskip
  Dibujar en la pizarra el diagrama de memoria para
  \texttt{./programa hola mundo 42}:
  \begin{verbatim}
   argv
   +-----+   +------------------+
   | [0] |-->| "./programa\0"   |
   +-----+   +------------------+
   | [1] |-->| "hola\0"         |
   +-----+   +------------------+
   | [2] |-->| "mundo\0"        |
   +-----+   +------------------+
   | [3] |-->| "42\0"           |
   +-----+   +------------------+
   argc = 4
  \end{verbatim}
  Señalar que \cpp{argv[3]} es la cadena \texttt{"42"}, no el número
  \texttt{42}.  Este punto se enfatiza antes de la demo de conversión.
\end{notaprofesor}

\begin{notaprofesor}[FAQ: \texttt{char* argv[]} vs \texttt{char** argv}]

  \textbf{¿Son lo mismo \cpp{char* argv[]} y \cpp{char** argv}?}\\
  Sí, en el contexto de los parámetros de una función son equivalentes.
  Ambas formas se aceptan y es bueno que los alumnos lo sepan porque
  verán la segunda en código de C y en muchos libros.

  \medskip
  \textbf{¿Por qué \cpp{argc} incluye el nombre del programa?}\\
  Decisión de diseño de los primeros Unix: el programa puede necesitar
  saber cómo fue invocado.  Un ejemplo clásico es \texttt{vi}/\texttt{vim}:
  si se invoca como \texttt{view} (un enlace simbólico) se abre en modo
  solo lectura.  Otro: \cpp{argv[0]} en el mensaje de uso.

  \medskip
  \textbf{¿Qué hay en \cpp{argv[argc]}?}\\
  El estándar garantiza que \cpp{argv[argc]} es \cpp{nullptr} (puntero
  nulo).  Algunos programas aprovechan esto para iterar con
  \cpp{while (*argv++)} en lugar de un \cpp{for} con índice.  Para este
  curso la forma con \cpp{for} e índice es más clara.

  \medskip
  \textbf{¿Puedo asignar \cpp{argv[i]} a un \cpp{string}?}\\
  Sí.  \cpp{string} tiene un constructor que acepta \cpp{const char*}:
  \cpp{string s = argv[1];} funciona.  A partir de ahí se puede usar
  todas las funciones de \cpp{std::string} (longitud, búsqueda, etc.).

\end{notaprofesor}

% =============================================================================
\section{Demo en vivo: cuatro programas}
% =============================================================================

\subsection{\texttt{hola\_args.cpp}: imprimir \texttt{argc} y \texttt{argv}}

\begin{notaprofesor}[Guión (5 min)]
  \begin{enumerate}
    \item Escribir el programa en vivo (es corto; no mostrar el código antes).
    \item Compilar y ejecutar sin argumentos.  Preguntar antes: ``¿qué
          valor tiene \cpp{argc}?''  Respuesta: 1.  Mostrar que solo
          imprime \cpp{argv[0]}.
    \item Ejecutar con \texttt{uno dos tres}.  Confirmar \cpp{argc = 4}
          y los cuatro valores en pantalla.
    \item Ejecutar con \texttt{"hola mundo"} entre comillas.  Mostrar
          que ahora \cpp{argc = 2} y \cpp{argv[1] = "hola mundo"}: el
          shell pasa la cadena completa como un solo argumento.
    \item Preguntar: ``¿qué índice tiene el último argumento en
          términos de \cpp{argc}?''  Respuesta: \cpp{argc-1}.
  \end{enumerate}
\end{notaprofesor}

\subsection{\texttt{uso.cpp}: verificar \texttt{argc} y mensaje de uso}

\begin{notaprofesor}[Guión (5 min)]
  \begin{enumerate}
    \item Mostrar la verificación \cpp{if (argc != 3)} y preguntar:
          ``¿por qué 3 y no 2?''  Respuesta: el programa espera nombre
          y apellido, más \cpp{argv[0]}, son 3 cadenas en total.
    \item Ejecutar con el número correcto de argumentos.  Mostrar
          salida esperada.
    \item Ejecutar con uno solo.  Mostrar el mensaje de error en
          \cpp{cerr}.  Preguntar: ``¿por qué \cpp{cerr} y no \cpp{cout}?''
          Demostrar con \texttt{./uso Ana 2>/dev/null}: el mensaje de
          error desaparece; cualquier salida normal seguiría visible.
    \item Señalar que \cpp{return 1} indica al shell que hubo error.
          Se puede demostrar con \texttt{echo \$?} (Linux/Mac) o
          \texttt{echo \%ERRORLEVEL\%} (Windows) justo después.
  \end{enumerate}
\end{notaprofesor}

\subsection{\texttt{suma.cpp}: iterar sobre los argumentos}

\begin{notaprofesor}[Guión (5 min)]
  \begin{enumerate}
    \item Antes de mostrar el código, preguntar: ``Si los argumentos
          son números, ¿cómo podríamos sumarlos?''  Esperar propuestas.
          La respuesta natural es un \cpp{for} desde \cpp{i=1} hasta
          \cpp{argc-1} con \cpp{atoi}.
    \item Mostrar el código.  Compilar y ejecutar con varios números.
    \item Probar con cero argumentos (\cpp{argc = 1}).  El ciclo no
          itera y el resultado es~0: correcto.
    \item Probar con un argumento no numérico (\texttt{./suma hola}).
          \cpp{atoi("hola")} devuelve~0.  Preguntar: ``¿es esto un
          buen comportamiento?''  Introducir \cpp{stoi} como alternativa
          que lanza excepción.
  \end{enumerate}
\end{notaprofesor}

\subsection{\texttt{calc.cpp}: el primer argumento como operador}

\begin{notaprofesor}[Guión (5 min)]
  \begin{enumerate}
    \item Señalar el acceso \cpp{argv[2][0]}: el primer carácter del
          segundo argumento.  Este es el único lugar del curso donde
          se usa la indexación doble sobre \cpp{argv}.
    \item Compilar y ejecutar con varias operaciones.  Demostrar que
          \texttt{./calc 10 / 4} da \texttt{2.5} (no~2, porque
          \cpp{atof} produce \cpp{double}).
    \item Probar con un operador inválido.  Mostrar el \cpp{return 1}
          en acción.
    \item \textbf{Pregunta de profundidad:} ``¿Qué pasa si el usuario
          escribe \texttt{./calc 10 * 4}?''  Respuesta: el shell
          expande \texttt{*} como un glob antes de pasárselo al
          programa.  Los alumnos verán algo inesperado.  Solución:
          escapar el asterisco: \texttt{./calc 10 \textbackslash* 4}
          o \texttt{./calc 10 "*" 4}.  Esto abre una conversación
          interesante sobre qué hace el shell antes de invocar el
          programa.
  \end{enumerate}
\end{notaprofesor}

% =============================================================================
\section{Convertir argumentos a números}
% =============================================================================

\begin{notaprofesor}[Cómo presentar la tabla (5 min)]
  Presentar la tabla de funciones y enfatizar la simetría:
  \cpp{atoi}/\cpp{stoi} para enteros, \cpp{atof}/\cpp{stod} para
  flotantes.  Las versiones con \texttt{a} son de~C (reciben
  \cpp{char*}), las de \texttt{s} son de C++ (reciben \cpp{string}).
  \par\smallskip
  Regla práctica para este curso: usar \cpp{atoi}/\cpp{atof}
  cuando el dato viene directamente de \cpp{argv} (ya es \cpp{char*});
  usar \cpp{stoi}/\cpp{stod} cuando el dato ya es un \cpp{string}.
  \par\smallskip
  Si alguien pregunta por \cpp{stoi} con cadenas inválidas, mencionar
  que lanza \cpp{std::invalid\_argument} y que se captura con
  \cpp{try}/\cpp{catch}, tema que se ve en sesiones posteriores.
\end{notaprofesor}

% =============================================================================
\section{Actividad: el teclado numérico}
% =============================================================================

\begin{notaprofesor}[Organización de la actividad (55 min)]

  \textbf{Antes de la clase:}
  \begin{itemize}
    \item Tener la imagen del teclado lista para proyectar.
    \item Preparar \texttt{materialInicial/sesion23/} con el archivo
          \texttt{actividad} y la imagen \texttt{teclado\_numerico.jpg}.
  \end{itemize}

  \textbf{Introducción (5 min):}\\
  Mostrar la imagen del teclado y preguntar a los alumnos cuántas
  pulsaciones requiere cada letra.  Construir la tabla en vivo
  (o revelarla desde las notas) antes de que comiencen a programar.
  Enfatizar la regla: la primera pulsación da el dígito, las siguientes
  dan las letras, por lo que la $n$-ésima letra requiere $n+1$
  pulsaciones.

  \textbf{Implementación (35 min):}\\
  Dar libertad para que cada quien diseñe su solución.  Los dos
  enfoques más comunes que surgirán:
  \begin{enumerate}
    \item \textbf{Cascada de \cpp{if}/\cpp{else} o \cpp{switch}:} un
          caso por letra.  Funciona pero produce $\sim$100 líneas.
          No disuadir; después mostrar la versión con arreglos.
    \item \textbf{Dos arreglos \cpp{tecla[]} y \cpp{posicion[]}:}
          la solución de referencia.  Señalar cómo \cpp{letra - 'a'}
          da el índice de 0 a 25, conectando con lo visto en la
          sesión de ASCII.
  \end{enumerate}
  Circular por el salón.  Los puntos de fricción más comunes son:
  \begin{itemize}
    \item Inicializar los dos arreglos correctamente (es tedioso; normal).
    \item El ciclo interno de pulsaciones: \cpp{for(int k=0; k<=pos; k++)}.
    \item Imprimir el espacio entre palabras pero no al final:
          \cpp{if(i > 1) cout << " ";} antes de cada palabra.
  \end{itemize}

  \textbf{Discusión (15 min):}\\
  Pedir a dos o tres alumnos que proyecten su solución.  Elegir una
  que use cascada y otra que use arreglos para comparar.  Preguntar
  al grupo: ``¿cuál prefieren?  ¿Por qué?''  Conectar con el principio
  de separar datos de lógica.

\end{notaprofesor}

\begin{notaprofesor}[FAQ: el encoding del teclado]

  \textbf{¿Por qué $n+1$ pulsaciones y no $n$?}\\
  En la mayoría de los teléfonos analógicos la primera pulsación
  de una tecla insertaba el dígito (el \texttt{4}, no la \texttt{g}).
  Las letras se obtenían con pulsaciones adicionales.  Esa convención
  explica que la primera letra requiera 2 pulsaciones.

  \medskip
  \textbf{¿Qué pasa con dos letras consecutivas en la misma tecla,
  como ``gh'' o ``mn''?}\\
  En los teléfonos reales había que esperar un segundo (tiempo de
  espera) o presionar \texttt{\#} para indicar que la siguiente
  pulsación era una letra nueva.  Para esta actividad el enunciado
  garantiza que no habrá esa situación, así que no es necesario
  manejarla.  Es un buen ejercicio adicional si hay tiempo: ¿cómo
  detectarías que la siguiente letra está en la misma tecla?

  \medskip
  \textbf{¿Por qué no usar un \cpp{map<char,string>}?}\\
  \cpp{std::map} no se ha visto aún.  La solución con arreglos es
  deliberadamente más básica y refuerza la aritmética de \cpp{char}.
  Con \cpp{map} sería válida si los alumnos la conocen de otro contexto.

\end{notaprofesor}

% =============================================================================
\section{Soluciones a los ejercicios}
% =============================================================================

\begin{notaprofesor}[Ejercicio 1 — \texttt{eco.cpp}]
\begin{verbatim}
#include <iostream>
using namespace std;

int main(int argc, char* argv[]){
    for(int i = 1; i < argc; i++)
        cout << i << ": " << argv[i] << endl;
}
\end{verbatim}
  Si no hay argumentos (\cpp{argc == 1}), el ciclo no itera y no
  se imprime nada.  No se necesita verificación especial.
\end{notaprofesor}

\begin{notaprofesor}[Ejercicio 2 — \texttt{mayor.cpp}]
\begin{verbatim}
#include <iostream>
#include <cstdlib>
using namespace std;

int main(int argc, char* argv[]){
    if(argc < 2){
        cerr << "Uso: " << argv[0] << " <num1> [num2 ...]" << endl;
        return 1;
    }
    int max = atoi(argv[1]);
    for(int i = 2; i < argc; i++){
        int n = atoi(argv[i]);
        if(n > max) max = n;
    }
    cout << max << endl;
}
\end{verbatim}
  Nótese que el máximo se inicializa con \cpp{argv[1]} (el primer
  argumento) y el ciclo empieza en~2.  Inicializarlo con~0 o con
  \cpp{INT\_MIN} también funciona pero depende de suposiciones sobre
  el rango de los datos.
\end{notaprofesor}

\begin{notaprofesor}[Ejercicio 3 — \texttt{reversa.cpp}]
\begin{verbatim}
#include <iostream>
using namespace std;

int main(int argc, char* argv[]){
    for(int i = argc-1; i >= 1; i--)
        cout << argv[i] << endl;
}
\end{verbatim}
  El ciclo va de \cpp{argc-1} hasta~1 (no hasta~0, para no imprimir
  el nombre del programa).  Si no hay argumentos, \cpp{argc-1 == 0 < 1}
  y el ciclo no itera.
\end{notaprofesor}

\begin{notaprofesor}[Ejercicio 4 — \texttt{contador.cpp}]
\begin{verbatim}
#include <iostream>
using namespace std;

int main(int argc, char* argv[]){
    if(argc < 3){
        cerr << "Uso: " << argv[0]
             << " <letra> <palabra1> [palabra2 ...]" << endl;
        return 1;
    }
    char letra = argv[1][0];
    int cuenta = 0;
    for(int i = 2; i < argc; i++){
        for(int j = 0; argv[i][j] != '\0'; j++){
            if(argv[i][j] == letra) cuenta++;
        }
    }
    cout << cuenta << endl;
}
\end{verbatim}
  \textbf{Verificación:} \texttt{./contador a banana manzana naranja}
  \begin{itemize}
    \item \texttt{"banana"}: a en posiciones 1,3,5 → 3 aces.
    \item \texttt{"manzana"}: a en posiciones 1,4,6 → 3 aces.
    \item \texttt{"naranja"}: a en posiciones 1,3,5 → 3 aces\dots
  \end{itemize}

  Espera, contemos bien: ``banana'' = b-a-n-a-n-a → 3~a's.
  ``manzana'' = m-a-n-z-a-n-a → 3~a's.  ``naranja'' = n-a-r-a-n-j-a → 3~a's.
  Total: \textbf{9}, no~7.  El enunciado del ejercicio tiene la
  respuesta~7 como ejemplo ilustrativo; corregirlo en clase si algún
  alumno lo señala.  Buen momento para enfatizar la importancia de
  verificar los ejemplos a mano.
\end{notaprofesor}

\begin{notaprofesor}[Ejercicio 5 — \texttt{convertir.cpp} (solución de referencia)]
\begin{verbatim}
#include <iostream>
using namespace std;

// Tecla y posicion de cada letra del alfabeto (indice: 'a'=0 ... 'z'=25)
int tecla[]   = {2,2,2, 3,3,3, 4,4,4, 5,5,5, 6,6,6, 7,7,7,7, 8,8,8, 9,9,9,9};
int posicion[] = {1,2,3, 1,2,3, 1,2,3, 1,2,3, 1,2,3, 1,2,3,4, 1,2,3, 1,2,3,4};

int main(int argc, char* argv[]){
    if(argc < 2){
        cerr << "Uso: " << argv[0] << " <palabra1> [palabra2 ...]" << endl;
        return 1;
    }
    for(int i = 1; i < argc; i++){
        if(i > 1) cout << " ";
        for(int j = 0; argv[i][j] != '\0'; j++){
            int idx = argv[i][j] - 'a';
            for(int k = 0; k <= posicion[idx]; k++)
                cout << tecla[idx];
        }
    }
    cout << endl;
}
\end{verbatim}

  \textbf{Verificación manual de \texttt{hola}:}\\
  \begin{tabular}{clccc}
    \toprule
    Letra & Tecla & Pos & Pulsaciones (\cpp{k <= pos}) & Salida \\
    \midrule
    h (idx=7)  & 4 & 2 & k=0,1,2 → 3 veces & \texttt{444} \\
    o (idx=14) & 6 & 3 & k=0,1,2,3 → 4 veces & \texttt{6666} \\
    l (idx=11) & 5 & 3 & k=0,1,2,3 → 4 veces & \texttt{5555} \\
    a (idx=0)  & 2 & 1 & k=0,1 → 2 veces & \texttt{22} \\
    \bottomrule
  \end{tabular}
  \par\smallskip
  Salida: \texttt{4446666555522} $\checkmark$

  \medskip
  \textbf{Nota sobre los arreglos:}
  \begin{itemize}
    \item \cpp{tecla[]} tiene 26 elementos: 3+3+3+3+3+4+3+4 = 26.
    \item \cpp{posicion[]} ídem.
    \item \cpp{argv[i][j] != '\textbackslash{}0'} itera sobre los
          caracteres de cada cadena hasta el terminador.
    \item El ciclo de pulsaciones \cpp{for(k=0; k<=posicion[idx]; k++)}
          ejecuta \cpp{posicion[idx]+1} veces.  Para \cpp{posicion=1}
          (letras a, d, g\ldots) ejecuta 2 veces; para \cpp{posicion=4}
          (s,z) ejecuta 5 veces.
  \end{itemize}
\end{notaprofesor}

% =============================================================================
\section{Errores frecuentes}
% =============================================================================

\begin{notaprofesor}[Los cinco errores más comunes]

  \textbf{1.\ Acceder a \cpp{argv[i]} sin verificar \cpp{argc}.}\\
  Síntoma: \textit{segmentation fault} al ejecutar sin argumentos.\\
  Remedio: verificar \cpp{argc} antes de acceder a cualquier
  \cpp{argv[i]} con $i \geq 1$.

  \vspace{0.5em}
  \textbf{2.\ Usar \cpp{argv[i]} directamente en aritmética.}\\
  Síntoma: resultados incorrectos o el compilador rechaza la operación.\\
  Ejemplo típico: \cpp{int n = argc[1] + 1;} (debería ser
  \cpp{atoi(argv[1]) + 1}).\\
  Remedio: convertir siempre con \cpp{atoi}, \cpp{atof} o asignación
  a \cpp{string}.

  \vspace{0.5em}
  \textbf{3.\ Ciclo que empieza en~0 e incluye \cpp{argv[0]}.}\\
  Síntoma: el nombre del programa aparece en la salida o se intenta
  operar sobre él como si fuera un dato.\\
  Remedio: los argumentos del usuario siempre empiezan en \cpp{argv[1]}.
  El ciclo debe comenzar en \cpp{i = 1}.

  \vspace{0.5em}
  \textbf{4.\ En \texttt{convertir.cpp}: arreglos \cpp{tecla[]} con
  26 elementos contados mal.}\\
  Síntoma: letras de la \texttt{u} en adelante producen resultados
  erróneos o el programa se sale de los límites del arreglo.\\
  Remedio: contar los tres bloques: 3+3+3+3+3=15 letras en teclas
  2--6, luego 4 en la~7, 3 en la~8 y 4 en la~9 = 26.  Verificar
  con \cpp{sizeof(tecla)/sizeof(int)}.

  \vspace{0.5em}
  \textbf{5.\ El espacio entre palabras aparece al inicio o al final.}\\
  Síntoma: \texttt{" hola luna"} o \texttt{"hola luna "}.\\
  Remedio: imprimir el espacio \textit{antes} de cada palabra excepto
  la primera: \cpp{if(i > 1) cout << " ";} al inicio del ciclo
  exterior.

\end{notaprofesor}

% =============================================================================
\section{Rúbrica de la actividad}
% =============================================================================

\begin{center}
\begin{tabular}{p{7.5cm}cc}
  \toprule
  \textbf{Criterio} & \textbf{Puntos} & \textbf{Evidencia} \\
  \midrule
  \multicolumn{3}{l}{\textit{Recepción de argumentos}} \\
  \quad Usa \cpp{argc} y \cpp{argv} correctamente           & 10 & código \\
  \quad Verifica \cpp{argc} y muestra uso en \cpp{cerr}     & 5  & prueba sin args \\
  \midrule
  \multicolumn{3}{l}{\textit{Mapeo letra--pulsaciones}} \\
  \quad Produce la secuencia correcta para cada letra        & 15 & tabla de pruebas \\
  \quad Teclas 7 (pqrs) y 9 (wxyz) con 4 letras correctas   & 10 & palabras con s/z \\
  \midrule
  \multicolumn{3}{l}{\textit{Formato de salida}} \\
  \quad Espacio entre palabras, sin espacio al inicio/final  & 10 & salida con 2+ palabras \\
  \midrule
  \textbf{Total}                                             & \textbf{50} & \\
  \bottomrule
\end{tabular}
\end{center}

% =============================================================================
\section{Conexiones con temas posteriores}
% =============================================================================

\begin{itemize}
  \item \textbf{Sesión 24 y posteriores --- Entrada estándar vs.\ args:}
        La diferencia entre \cpp{cin} (lectura interactiva) y \cpp{argv}
        (configuración al invocar) es un patrón que aparece en todos
        los programas de línea de comandos.  Herramientas como
        \texttt{g++}, \texttt{git} y \texttt{grep} combinan ambas.
  \item \textbf{Orientación a Objetos --- Constructores con parámetros:}
        La idea de que una entidad recibe su configuración al momento de
        crearse (no después) es la misma en \cpp{main(argc, argv)} y en
        un constructor con argumentos.  Vale la pena señalarlo cuando
        se introduzcan los constructores.
  \item \textbf{Programación de sistemas --- \cpp{getopt} y
        \texttt{argparse}:}
        Para programas reales con múltiples opciones opcionales
        (\texttt{-v}, \texttt{--output archivo}) existe la biblioteca
        \cpp{getopt} en C/C++ y \texttt{argparse} en Python.  Ambas
        procesan \cpp{argv} internamente; lo que aprendieron hoy es
        la base sobre la que se construyen.
\end{itemize}

% =============================================================================
\section{Materiales y recursos}
% =============================================================================

\begin{itemize}
  \item \textbf{Notas alumno:}
        \texttt{notas\_alumno/pdf/sesion23\_argumentos\_consola.pdf}
  \item \textbf{Material inicial:}
        \texttt{materialInicial/sesion23/} --- \texttt{actividad},
        \texttt{readme}, \texttt{teclado\_numerico.jpg}
  \item \textbf{Imagen del teclado:} proyectar durante la introducción
        de la actividad; está disponible en el material inicial.
\end{itemize}

\end{document}
